% ISEF Poster — LaTeX template using tikzposter
% Build: xelatex isef-poster-template.tex
% (xelatex preferred for Unicode/Chinese; pdflatex also works for English-only)
%
% Adjust paper size in the \documentclass line:
%   [a0paper,portrait] for A0 portrait (841×1189mm) — DEFAULT
%   [a0paper,landscape] for A0 landscape (1189×841mm)
%   Custom: use \setlength{\paperwidth}{914mm}\setlength{\paperheight}{1219mm} for 36"×48"
%
% Fits ISEF booth (122×165cm practical max) as long as you stay ≤A0 portrait or 48×36" landscape.

\documentclass[a0paper,portrait,18pt]{tikzposter}

\usepackage[utf8]{inputenc}
\usepackage[T1]{fontenc}
\usepackage{lmodern}        % cleaner fonts
\usepackage{xeCJK}          % comment out if not using CJK
\setCJKmainfont{PingFang SC}[AutoFakeBold=true]

\usepackage{amsmath,amssymb}
\usepackage{graphicx}
\usepackage{hyperref}
\usepackage{qrcode}
\usepackage{xcolor}

% --- ISEF color theme (navy + orange accent; meets WCAG 4.5:1 on white) ---
\definecolor{isefNavy}{HTML}{0B3D91}
\definecolor{isefOrange}{HTML}{FFA500}
\definecolor{isefGray}{HTML}{4A4A4A}

% Title formatting — large title, readable from 3m
\settitle{
  \centering
  \color{white}{\fontsize{100pt}{120pt}\selectfont\textbf{\@title}}\\[1cm]
  \color{white}{\fontsize{36pt}{42pt}\selectfont \@author}\\[0.3cm]
  \color{white}{\fontsize{28pt}{32pt}\selectfont \textit{\institute}}
}

% Apply theme
\definecolorpalette{ISEF}{
  \definecolor{colorOne}{named}{isefNavy}
  \definecolor{colorTwo}{named}{isefOrange}
  \definecolor{colorThree}{named}{white}
}
\usecolorpalette{ISEF}
\usetheme{Default}
\useblockstyle{Default}

% --- Project metadata ---
\title{Origami Configurations as Statistical Ensembles: A Free-Energy Analysis}
\author{Hikaru Kuribayashi}
\institute{Stuyvesant High School • PHYS — Physics and Astronomy • PHYS021}

\begin{document}

\maketitle[width=0.96\textwidth]

% ---- Top row: abstract + significance ----
\begin{columns}
  \column{0.65}
  \block{Abstract}{
    \fontsize{24pt}{30pt}\selectfont
    Origami crease patterns admit many physical configurations of the folded sheet.
    We recast this configurational ensemble as a statistical-mechanics partition
    function and use Markov-chain Monte Carlo (MCMC) to sample states. The
    Metropolis sampler reproduces known closed-form distributions for 2-fold
    linkages (validation set), then extends to multi-fold patterns without
    closed-form solutions. We find a discrete-to-continuous phase transition at
    fold density $\rho \approx 0.4$ that has not been previously reported in
    the origami literature.
  }
  \column{0.35}
  \block{Significance}{
    \fontsize{24pt}{30pt}\selectfont
    This work imports statistical mechanics into origami, opening connections
    to crystallization theory, polymer self-assembly, and protein folding.
    Code and data are open-source.

    \vspace{0.5cm}
    \centering
    \qrcode[height=4cm]{https://github.com/example/origami-mcmc}\\
    \vspace{0.2cm}
    {\fontsize{20pt}{24pt}\selectfont \textit{Code + data}}
  }
\end{columns}

% ---- Three-column body ----
\begin{columns}
  \column{0.33}
  \block{Background}{
    \fontsize{22pt}{28pt}\selectfont
    Origami's mathematics has been studied for foldability (Demaine 2007) and
    rigidity, but the space of accessible folded states has rarely been treated
    as a probabilistic ensemble. Statistical mechanics provides ready tools:
    partition functions, free energies, phase transitions.
  }
  \block{Methods}{
    \fontsize{22pt}{28pt}\selectfont
    Each fold = continuous angle $\theta_i \in [0,\pi]$.
    Hamiltonian $H(\boldsymbol{\theta}) = \sum_i \kappa_i \theta_i^2 + V_\text{constraint}$
    where $V_\text{constraint}$ penalizes infeasible folds (self-intersection,
    crease violations).

    \vspace{0.3cm}
    MCMC: Metropolis-Hastings with adaptive step size; 10$^6$ samples after
    burn-in of 10$^4$; thinning at $\tau \approx 50$.
  }

  \column{0.34}
  \block{Validation: 2-fold linkage}{
    \fontsize{22pt}{28pt}\selectfont
    \centering
    \fbox{\parbox{0.95\linewidth}{\centering\textit{[Place \texttt{figures/validation-2fold.png} here]}}}\\
    \vspace{0.3cm}
    Sampled distribution (orange) matches the closed-form distribution (blue)
    derived from Lechenault et al.\ (2014).
  }
  \block{Phase transition at $\rho \approx 0.4$}{
    \fontsize{22pt}{28pt}\selectfont
    \centering
    \fbox{\parbox{0.95\linewidth}{\centering\textit{[Place \texttt{figures/phase-transition.png} here]}}}\\
    \vspace{0.3cm}
    Free energy as a function of fold density. Discontinuity at
    $\rho_c \approx 0.4$ corresponds to topological reorganization of the
    configuration space.
  }

  \column{0.33}
  \block{Discussion}{
    \fontsize{22pt}{28pt}\selectfont
    The discrete-to-continuous transition arises from the topology of the
    configuration space. Below $\rho_c$ only discrete fold patterns are
    accessible; above $\rho_c$ a continuum opens up. This may relate to the
    crystalline$\rightarrow$amorphous transition in folded polymer self-assembly.
  }
  \block{Future work}{
    \fontsize{22pt}{28pt}\selectfont
    \begin{itemize}\setlength\itemsep{0.3em}
      \item 3D fold patterns (current work is 2D-flat)
      \item Connection to protein-folding energy landscapes
      \item Experimental validation via micro-origami
    \end{itemize}
  }
\end{columns}

% ---- Footer ----
\block{References \& Acknowledgments}{
  \fontsize{16pt}{20pt}\selectfont
  \begin{minipage}[t]{0.55\linewidth}
  \textbf{References}
  \begin{enumerate}\setlength\itemsep{0.1em}
    \item Demaine, E. \& O'Rourke, J. (2007). \textit{Geometric Folding Algorithms}. Cambridge UP.
    \item Tachi, T. (2010). Origami simulator. \textit{PRSA} 466.
    \item Newman, M. \& Barkema, G. (1999). \textit{Monte Carlo Methods in Statistical Physics}. OUP.
    \item Lechenault, F.\ et al.\ (2014). Mechanical response of a creased sheet. \textit{PRL} 112.
    \item Silverberg, J.\ et al.\ (2014). Using origami design principles. \textit{Science} 345.
  \end{enumerate}
  \end{minipage}
  \hfill
  \begin{minipage}[t]{0.40\linewidth}
  \textbf{Acknowledgments}\\
  Thanks to my advisor Dr.\ R.\ Smith for early feedback. MCMC code is original;
  validation against Demaine et al.\ closed forms confirms correctness.

  \vspace{0.3cm}
  \textbf{ISEF 2026 AI-use disclosure:}
  Claude assisted with code review and manuscript polish. Methodology design,
  theoretical framing, experiments, and interpretation are my own work.
  \end{minipage}
}

\end{document}
